% ===================================================================
%  QUADRO N.º 5 — A Casa e a sua construção.
%
%  Ceci n'est PAS une transcription : c'est une traduction, faite en
%  2026, en portugais d'aujourd'hui. Voir texto/pt/10-quadro-01.tex
%  pour la consigne, qui vaut pour les seize fichiers.
%
%  LES DEUX NOMS SONT DES COULEURS, en francais comme partout :
%  M. Leblanc et M. Roux sont ici « senhor Branco » et « senhor
%  Ruivo » — deux patronymes portugais reels, et qui disent la meme
%  chose que les deux autres.
% ===================================================================

%%K t05-apar-1 apar
\VUsaut{6.0mm}
\VUtitre{15.0pt}{60}{QUADRO N.º 5}
\VUsaut{1.4mm}
\VUtitre{11.4pt}{0}{A Casa e a sua construção.}
\VUsaut{2.2mm}

%%K t05-tit-1 sub
\VUpk{10.2pt}{40}{Um bom encontro.}
\VUsaut{3.0mm}

%%K t05-01-1 p
\textsc{João}. --- Tio, gostaria muito de saber como se constrói uma
\VUgras{casa}.

%%K t05-01-2 p
\textsc{O Tio} \textsuperscript{(80)}. --- É muito fácil, meu filho; vem comigo
e mostrar-te-ei a que o senhor Bernardo \textsuperscript{(79)} está a mandar
construir e na qual vou morar.

%%K t05-01-3 p
\textsc{João}. --- E quem desenhou a \VUgras{planta} \textsuperscript{(76)} desta
casa?

%%K t05-01-4 p
\textsc{O Tio}. --- O \VUgras{arquiteto} \textsuperscript{(75)}; apresentou um
desenho muito claro, que foi aprovado, e o \VUgras{empreiteiro}
\textsuperscript{(77)} começou as obras.

%%K t05-01-5 p
\textsc{João}. --- E quem é o empreiteiro?

%%K t05-01-6 p
\textsc{O Tio}. --- É o senhor Branco, o pai do teu amigo Jaime. Dirige os
operários com o senhor Ruivo, o arquiteto.

%%K t05-01-7 p
Ora vejam! Aqui chega justamente o senhor Branco com a sua \VUgras{régua}
\textsuperscript{(78)}, e o senhor Ruivo com a sua planta.

%%K t05-01-8 p
Bom dia, senhores. Como estão?

%%K t05-01-9 p
\textsc{O Sr. Ruivo}. --- Muito bem, obrigado. Bom dia, João; como este
menino cresceu!

%%K t05-01-10 p
\textsc{O Tio}. --- Pudera, é todo um homenzinho. É muito sério; é preciso
explicar-lhe tudo o que vê. Hoje queria saber como se constrói uma casa.

%%K t05-tit-2 sub
\VUpk{10.2pt}{40}{Os operários.}
\VUsaut{3.0mm}

%%K t05-01-11 p
\textsc{O Sr. Branco}. --- Então vem comigo, meu amiguinho, e explicar-to-ei
já a seguir, porque gosto muito dos meninos que querem aprender.

%%K t05-01-12 p
Vês aquele operário que tem uma \VUgras{pá} \textsuperscript{(82)} na mão: é um
\VUgras{cantoneiro} \textsuperscript{(81)}. Está a abrir uma \VUgras{vala}
\textsuperscript{(46)} para colocar a canalização da água. Foi também ele que
cavou o solo no sítio onde está a casa, para que o pedreiro levantasse as
quatro \VUgras{paredes}. A parte dessas paredes que fica debaixo da terra
chama-se os \VUgras{alicerces} \textsuperscript{(2)}. O espaço compreendido entre
os alicerces forma a \VUgras{adega} \textsuperscript{(3)}. Essa adega tem uma
janelinha chamada \VUgras{respiradouro} \textsuperscript{(5)}, uma \VUgras{porta}
\textsuperscript{(4)} e uma escada.

%%K t05-01-13 p
O \VUgras{pedreiro} \textsuperscript{(108)} continuou a levantar as paredes
deixando vãos para as \VUgras{portas} \textsuperscript{(8)} e as \VUgras{janelas}
\textsuperscript{(17)}.

%%K t05-01-14 p
\textsc{João}. --- Mas não o vejo, a esse pedreiro!

%%K t05-01-15 p
\textsc{O Sr. Branco}. --- Já acabou de trabalhar na casa. Está junto à
\VUgras{cavalariça} \textsuperscript{(54)}, sobre o seu \VUgras{andaime}
\textsuperscript{(110)}, a levantar a parede do \VUgras{lavadouro}
\textsuperscript{(56)}.

%%K t05-01-16 p
Tem uma talocha, uma masseira e um \VUgras{prumo} \textsuperscript{(109)}. Mas não
te lembrarás desses nomes.

%%K t05-01-17 p
\textsc{João}. --- Lembrar-me-ei, sim, porque vou pedir ao meu tio uma
talocha pequena e uma masseira, e eu mesmo farei um prumo com um cordel.

%%K t05-tit-3 sub
\VUsaut{3.0mm}
\VUpk{10.2pt}{40}{Os materiais.}
\VUsaut{3.0mm}

%%K t05-01-18 p
\textsc{O Tio}. --- Ah, muito bem! Mas diz-me, João, o que é preciso para
construir uma casa?

%%K t05-01-19 p
\textsc{João}. --- São precisas \VUgras{pedras de cantaria} \textsuperscript{(59)} e
\VUgras{argamassa} \textsuperscript{(65)}.

%%K t05-01-20 p
\textsc{O Tio}. --- Exato. Olha para aquele homem do chapéu grande, no seu
\VUgras{carro} \textsuperscript{(136)}. É um \VUgras{carreteiro}
\textsuperscript{(135)}. Todos os dias traz aqui \VUgras{blocos de pedra}
\textsuperscript{(60)}. Descarrega o seu carro no pátio, e os \VUgras{canteiros}
\textsuperscript{(83)} lavram os blocos e serram-nos com a sua \VUgras{serra de
pedra} \textsuperscript{(85)}. Um \VUgras{servente} \textsuperscript{(146)} leva-os ao
pedreiro, que os coloca.

%%K t05-01-21 p
Entretanto, o \VUgras{amassador} \textsuperscript{(87)} prepara a argamassa.
Precisa de \VUgras{areia} \textsuperscript{(62)}, \VUgras{cal} \textsuperscript{(63)} e
\VUgras{água} \textsuperscript{(64)}. A cal está ao seu lado num \VUgras{saco}
\textsuperscript{(71)}; um \VUgras{ajudante de pedreiro} \textsuperscript{(111)}
traz-lhe a areia num \VUgras{carrinho de mão} \textsuperscript{(112)}, e o
\VUgras{aprendiz} \textsuperscript{(113)} está na bomba \textsuperscript{(58)}. Enche um
\VUgras{balde} \textsuperscript{(114)}. A argamassa estará pronta já a seguir. O
\VUgras{servente de gamela} \textsuperscript{(107)} pega nela e sobe-a pela escada
até ao andaime, onde um pedreiro acaba esse lado da casa.

%%K t05-01-22 p
E agora, como se chama o operário que faz o \VUgras{telhado}
\textsuperscript{(31)}?

%%K t05-01-23 p
\textsc{João}. --- É o \VUgras{telhador} \textsuperscript{(118)}.

%%K t05-tit-4 sub
\VUsaut{3.0mm}
\VUpk{10.2pt}{40}{O telhado da casa.}
\VUsaut{3.0mm}

%%K t05-01-24 p
\textsc{O Sr. Branco}. --- Sim, mas primeiro vem o \VUgras{carpinteiro}
\textsuperscript{(115)}.

%%K t05-01-25 p
O carpinteiro faz a \VUgras{armação} \textsuperscript{(35)}. Une as \VUgras{traves}
\textsuperscript{(37)} com \VUgras{cavilhas} \textsuperscript{(117)}. Aqueles operários
lá em cima, com os seus largos calções de veludo, são carpinteiros. Um
segura uma trave pequena e o outro corta uma cavilha com o seu
\VUgras{machado} \textsuperscript{(116)}. O que está junto à \VUgras{chaminé}
\textsuperscript{(41)} é o telhador. Prega ripas sobre os \VUgras{barrotes} (*), e
\VUgras{ardósias} \textsuperscript{(66)} sobre as ripas. Às vezes põe telhas.

%%K t05-noto-1 noto
\VUnotes{143.2mm}{%
(*) \VUgras{Balk-o} = alemão \textit{Balken}, inglês \textit{joist},
francês \textit{solive}, italiano \textit{travicello}, russo
\textit{balka}, espanhol e português \textit{viga}. Raiz técnica ainda
não adotada, mas que se propõe aqui para traduzir o sentido da palavra
francesa \textit{solive}, que não é um \textit{trabo} (francês
\textit{poutre}) nem uma trave pequena. É um madeiro esquadriado que
sustenta um soalho e um teto.%
}

%%K t05-01-26 p
O telhado da cavalariça é de \VUgras{telhas} \textsuperscript{(55)}, o da casa é
de \VUgras{ardósia} \textsuperscript{(32)} e o da galeria é de \VUgras{zinco}
\textsuperscript{(24)}.

%%K t05-01-27 p
\textsc{João}. --- O telhador faz também os telhados de zinco?

%%K t05-01-28 p
\textsc{O Sr. Branco}. --- Não, esse é o \VUgras{zingueiro} \textsuperscript{(121)}
ou o \VUgras{canalizador}, o que está a colocar uma \VUgras{mansarda}
\textsuperscript{(38)} no telhado da casa. Também põe as \VUgras{caleiras}
\textsuperscript{(43)} e os \VUgras{tubos de queda} \textsuperscript{(42)}. Solda o
zinco e a folha de flandres. Foi ele que fixou o \VUgras{para-raios}
\textsuperscript{(40)} e o \VUgras{cata-vento} \textsuperscript{(39)} na \VUgras{empena}
\textsuperscript{(36)}.

%%K t05-01-29 p
\textsc{João}. --- Então a casa está terminada?

%%K t05-tit-5 sub
\VUsaut{3.0mm}
\VUpk{10.2pt}{40}{As ferramentas.}
\VUsaut{3.0mm}

%%K t05-01-30 p
\textsc{O Sr. Branco}. --- Não de todo. O \VUgras{estucador}
\textsuperscript{(99)} levanta com \VUgras{tijolos} \textsuperscript{(61)} os tabiques
que a dividem nos diferentes quartos. Reveste de \VUgras{gesso}
\textsuperscript{(70)} os \VUgras{tetos} \textsuperscript{(26)} e as paredes. Agora está
a fazer o teto da \VUgras{galeria} \textsuperscript{(33)}. Tem, como o pedreiro,
uma \VUgras{talocha} \textsuperscript{(100)} e uma \VUgras{masseira}
\textsuperscript{(101)}.

%%K t05-01-31 p
Depois o marceneiro faz as portas, as janelas, os \VUgras{soalhos}
\textsuperscript{(16)} e as \VUgras{portadas} \textsuperscript{(19)}.

%%K t05-01-32 p
Agora trabalha no pátio na sua \VUgras{bancada de carpinteiro}
\textsuperscript{(90)}. Tem uma \VUgras{serra} \textsuperscript{(94)} para serrar a
\VUgras{madeira} \textsuperscript{(69)}, uma \VUgras{plaina} \textsuperscript{(91)}, um
\VUgras{grampo} \textsuperscript{(92)}, um \VUgras{formão} \textsuperscript{(94
\textit{bis})}, um \VUgras{compasso} \textsuperscript{(95)} e um \VUgras{maço de
madeira} \textsuperscript{(95 \textit{bis})}.

%%K t05-01-33 p
\textsc{João}. --- Ah, mas é ainda mais divertido ser marceneiro do que
pedreiro! Tio, comprar-me-ás uma bancada, uma serra e uma plaina? Vou
fazer uma casota para o Medor.

%%K t05-01-34 p
\textsc{O Tio}. --- Sim, meu filho.

%%K t05-01-35 p
\textsc{João}. --- Que contente estou! E aquele que está junto à porta de
entrada, quem é?

%%K t05-tit-6 sub
\VUsaut{3.0mm}
\VUpk{10.2pt}{40}{A pintura.}
\VUsaut{3.0mm}

%%K t05-01-36 p
\textsc{O Sr. Branco}. --- É um \VUgras{pintor} \textsuperscript{(102)}. Está na
sua escada com a sua lata de \VUgras{tinta} \textsuperscript{(103)} e a sua
\VUgras{trincha} \textsuperscript{(104)}. Pinta a madeira da porta de entrada
para que se conserve mais tempo.

%%K t05-01-37 p
É também ele que cola o papel nos quartos.

%%K t05-01-38 p
O \VUgras{estofador} \textsuperscript{(107)} que está numa \VUgras{escada}
\textsuperscript{(106)}, na \VUgras{varanda} \textsuperscript{(28)}, coloca um
\VUgras{estore} \textsuperscript{(29)}.

%%K t05-01-39 p
O operário que está junto à janela grande é o \VUgras{vidraceiro}
\textsuperscript{(96)}. Fixa os \VUgras{vidros} \textsuperscript{(18)} com
\VUgras{betume} \textsuperscript{(73)}. Aquele outro operário, ajoelhado junto à
porta da adega, é um \VUgras{serralheiro} \textsuperscript{(97)}. Está a pôr uma
\VUgras{fechadura} \textsuperscript{(6)}. A sua \VUgras{lima} \textsuperscript{(98)} está
ao seu lado.

%%K t05-01-40 p
\textsc{João}. --- É também serralheiro aquele que bate com o seu
\VUgras{martelo} \textsuperscript{(84)} sobre um ferro?

%%K t05-01-41 p
\textsc{O Sr. Branco}. --- É antes um \VUgras{ferreiro} \textsuperscript{(125)}.
Forja \VUgras{ferros} \textsuperscript{(67)} para o \VUgras{eletricista}
\textsuperscript{(124)}, que está no telhado da \VUgras{cocheira}
\textsuperscript{(51)}. Tem uma \VUgras{forja} \textsuperscript{(126)}, uma
\VUgras{bigorna} \textsuperscript{(127)} e umas \VUgras{tenazes}
\textsuperscript{(128)}.

%%K t05-01-42 p
O eletricista coloca o \VUgras{fio de cobre} \textsuperscript{(44)} do telefone.

%%K t05-tit-7 sub
\VUpk{10.2pt}{40}{A jardinagem.}
\VUsaut{3.0mm}

%%K t05-01-43 p
\textsc{O Tio}. --- Ao fundo do \VUgras{pátio} \textsuperscript{(45)}, junto à
\VUgras{grade} \textsuperscript{(47)} e ao \VUgras{portão} \textsuperscript{(48)}, vês o
\VUgras{jardineiro} \textsuperscript{(129)}. Está a preparar o pequeno
\VUgras{jardim} \textsuperscript{(49)}. Cavou a terra com a sua \VUgras{enxada de
bico} \textsuperscript{(131)}, semeia sementes e alisa com o \VUgras{ancinho}
\textsuperscript{(130)}. Depois, com o seu \VUgras{regador} \textsuperscript{(132)},
regará os \VUgras{canteiros} \textsuperscript{(150)} e as plantas crescerão.

%%K t05-01-44 p
O \VUgras{moço de cavalariça} \textsuperscript{(133)}, que acaba de tratar do
cavalo e de lhe dar de comer, limpa a \VUgras{carruagem} \textsuperscript{(52)}
com uma esponja, uma \VUgras{bomba de mão} \textsuperscript{(134)} e um balde de
água. Depois metê-la-á na cocheira e fechará a porta com o grosso
\VUgras{ferrolho} \textsuperscript{(53)}.

%%K t05-01-45 p
Já estás contente, suponho; tiveste explicações que chegue.

%%K t05-01-46 p
\textsc{João}. --- Sim, tio, mas gostaria de ver o interior da casa.

%%K t05-01-47 p
\textsc{O Tio}. --- Isso será amanhã. O senhor Ruivo explicar-te-á a
planta e dir-te-á o nome de cada quarto.

%%K t05-tit-8 sub
\VUsaut{3.0mm}
\VUpk{10.2pt}{40}{A distribuição interior da casa.}
\VUsaut{2.4mm}

%%K t05-apar-2 apar
{\centering\textit{(Ver a planta.)}\par}
\VUsaut{2.4mm}

%%K t05-01-48 p
\textsc{O Arquiteto}. --- Vem, João, vou fazer-te visitar as diferentes
partes da casa.

%%K t05-01-49 p
Entremos no \VUgras{rés do chão} \textsuperscript{(7)}. Aqui está o botão da
\VUgras{campainha} \textsuperscript{(11)}. Subimos os três \VUgras{degraus}
\textsuperscript{(14)} da \VUgras{escadaria} \textsuperscript{(9)}, passamos a
\VUgras{soleira} \textsuperscript{(10)} da \VUgras{porta de entrada}
\textsuperscript{(8)}, e ei-nos ao pé da \VUgras{escada} \textsuperscript{(13)}.

%%K t05-01-50 p
\textsc{João}. --- Isto é o corredor.

%%K t05-01-51 p
\textsc{O Arquiteto}. --- Não, isto é o \VUgras{vestíbulo} \textsuperscript{(12)}.
O \VUgras{corredor} \textsuperscript{(a)} está ao fundo. À direita ficam a
\VUgras{sala de visitas} \textsuperscript{(b)} e a \VUgras{sala de jantar}
\textsuperscript{(c)}; em frente da sala de jantar estão a \VUgras{cozinha}
\textsuperscript{(d)} e a \VUgras{copa} \textsuperscript{(e)}; depois, à frente, o
\VUgras{escritório} \textsuperscript{(f)}. A \VUgras{galeria} \textsuperscript{(23)}, com
a sua \VUgras{porta envidraçada} \textsuperscript{(25)}, está junto à sala de
visitas.

%%K t05-01-52 p
Agora subiremos a escada. Agarra-te ao \VUgras{corrimão}
\textsuperscript{(15)} e conta os degraus. São catorze. Aqui está o \VUgras{patamar}
\textsuperscript{(m)}: estamos no primeiro \VUgras{andar} \textsuperscript{(27)}.

%%K t05-tit-9 sub
\VUsaut{3.0mm}
\VUpk{10.2pt}{40}{O primeiro andar.}
\VUsaut{3.0mm}

%%K t05-01-53 p
Em frente da escada estão a \VUgras{retrete} \textsuperscript{(20)} e o
\VUgras{toucador} \textsuperscript{(j)}. À direita, um belo \VUgras{quarto de
dormir} \textsuperscript{(g)} com duas janelas para a \VUgras{fachada}
\textsuperscript{(1)} da casa. À esquerda estão a \VUgras{casa de banho}
\textsuperscript{(1)} e o \VUgras{quarto de hóspedes} \textsuperscript{(i)} com uma
grande \VUgras{alcova} \textsuperscript{(h)}. Junto ao quarto de dormir está o
\VUgras{quarto das crianças} \textsuperscript{(k)}. Dá para um amplo
\VUgras{terraço} \textsuperscript{(21)}. Por baixo desse terraço há outra
retrete \textsuperscript{(20)} e, contra a parede está o \VUgras{sino}
\textsuperscript{(22)} que chama para jantar.

%%K t05-01-54 p
\textsc{João}. --- Subamos ao \VUgras{segundo andar}.

%%K t05-01-55 p
\textsc{O Arquiteto}. --- Não há segundo andar. Sob o telhado só há
\VUgras{águas-furtadas} \textsuperscript{(33)} e o sótão \textsuperscript{(34)}.
Acabámos a visita, podemos descer.

%%K t05-01-56 p
\textsc{João}. --- Obrigado, senhor, foi muito interessante. Vou contar ao
meu pai tudo o que vi.
\VUsaut{4.0mm}
\VUfilet{22mm}
